%======================================================================%
\section{Sakharov-Induced Gravitation}
\label{sec:gravity:main}
%======================================================================%

Gravity in UCFT is not a fundamental input. It is the Sakharov
mechanism~\cite{Sakharov:1967pk}: integrating out the quantum
fluctuations of the matter sector over the curved emergent background
$M^{4}$ (the soldered 4-manifold of Postulate~P3, \autoref{sec:spacetime:main})
generates an Einstein--Hilbert term whose coefficient is the induced
Planck mass. The curvature dependence is carried entirely by the
second Seeley--DeWitt heat-kernel coefficient $a_{2}$, and the only
fields that feed the running $R$-coefficient are the \emph{massive}
ones. In UCFT the relevant massive multiplet is the set of
$28$ clock fields of common mass $m^{2}=f^{2}\mu^{2}$
(\autoref{sec:gravity:spectrum}), so the induced Newton constant is
fixed, positive, and computable. We work throughout in signature
$(-,+,+,+)$, with $[f]=\text{mass}$ and the dimensionless couplings
$\xi,\ \hat g^{2},\ \hat y^{2},\ \kappa,\ \Lambda_{\mathrm{cc}}$.

This section establishes three things: (i) the induced Planck mass
$M_{\mathrm{Pl}}^{2}=\tfrac{28}{6}\,f^{2}\mu^{2}\log(\Lambda^{2}/\mu^{2})/(16\pi^{2})>0$
(\emph{theorem}, given the spectrum and P3); (ii) that the
cosmological constant is \emph{not} solved --- the $a_{0}$ coefficient
forces an uncancelled $\sim\Lambda^{4}$ vacuum energy, handled only by a
tuned counterterm $\Lambda_{\mathrm{cc}}$; and (iii) the place of the
gravitational coupling $\kappa$ in the renormalization-group flow, whose
ultraviolet-attractive fixed point $\kappa_{\star}=96\pi^{2}/5$ is
conjectural. The full heat-kernel coefficient table and the
proper-time computation are deferred to \autoref{app:gravity:main}.

%----------------------------------------------------------------------%
\subsection{The massive multiplet and the ghost operator}
\label{sec:gravity:spectrum}
%----------------------------------------------------------------------%

Recall (\autoref{sec:potential:main}, Theorem~T4) that the vacuum manifold
is the rank-1 primitive-idempotent orbit
\[
  \mathcal M \;=\; \frac{E_{6}}{\mathrm{Spin}(10)\times U(1)} \;=\; \mathrm{EIII},
  \qquad \dim_{\mathbb R}\mathcal M = 32,
\]
with coset tangent module
$\mathfrak m = \mathbf{16}_{-3}\oplus\overline{\mathbf{16}}_{+3}$, a single
real-irreducible $H=\mathrm{Spin}(10)\times U(1)$ representation of complex
type (commutant $\mathbb C$), containing \emph{no} $\mathrm{Spin}(10)$
singlet (\autoref{sec:coset:main}). At tree level
$\Hess V_{\mathrm{tree}}|_{\mathfrak m}=0$: all $32$ clock fields are exact
Goldstones. The radiative spectrum is fixed by symmetry alone.

\begin{theorem}[Common radiative mass; $28+4$ split]
\label{thm:gravity:spectrum}
The one-loop Coleman--Weinberg effective potential is $H$-invariant.
Since $\mathfrak m$ is a single real-irreducible $H$-module, Schur's lemma
forces the radiative Hessian to be a multiple of the identity on all
$32$ directions,
\[
  \Hess V_{\mathrm{1\text{-}loop}}\big|_{\mathfrak m}
  \;=\; f^{2}\mu^{2}\,\mathbf 1_{32},
\]
i.e.\ a single common eigenvalue. Of these, four directions are the
frame/soldering modes that are gauge degrees of freedom (world-volume
diffeomorphism $+$ local Lorentz, Ward identity
$\nabla_{\mu}(\delta\Gamma/\delta e^{a}_{\mu})=0$); they are eaten by the
vierbein and local-Lorentz connection and carry \emph{no} physical mass.
The remaining
\[
  32 \;=\; 28\ \text{physical massive scalars (single common mass)}\;+\;4\ \text{gauge/soldering (eaten)}
\]
are the $28$ physical clock fields. With $[f]=\text{mass}$, the dim-4
Hessian eigenvalue is $f^{2}\mu^{2}$; the canonically normalized physical
mass is
\[
  m^{2} \;=\; \mu^{2}.
\]
The unbroken $H$ leaves the $46=\dim(\mathbf{45})+1$ gauge bosons massless.
\end{theorem}

\begin{proof}
$H$-invariance of $\Gamma$ and Schur's lemma on the irreducible $\mathfrak m$
give a scalar Hessian; this is the corrected version of the spectrum
analysis (\autoref{app:gravity:main}). The $4$ frame modes are not Hessian
zeros but gauge-protected directions, as established by the soldering Ward
identity of \autoref{sec:spacetime:main} (Theorem~3.1 there). Canonical
normalization divides the dim-4 eigenvalue $f^{2}\mu^{2}$ by $f^{2}$.
\end{proof}

\begin{remark}[No singlet subspace; against the legacy construction]
\label{rmk:gravity:nosinglet}
There is no four-dimensional $H$-invariant subspace of $\mathfrak m$ built
from $\nabla I_{2},\nabla I_{3}$ singlets: $\mathfrak m$ is irreducible and
has no $\mathrm{Spin}(10)$ singlet, so a singlet-projector construction is
algebraically impossible. The mass is one common eigenvalue from Schur, and
the four light modes are gauge-protected, not radiative zeros.
\end{remark}

The Euclidean fluctuation operator for the massive multiplet, used in the
heat-kernel expansion, is
\begin{equation}\label{eq:gravity:Deltam}
  \Delta_{m} \;=\; -D^{2} + m^{2}, \qquad m^{2}=f^{2}\mu^{2}
  \quad(\text{Hessian normalization}),
\end{equation}
with $D$ the $H$-covariant derivative on $\mathfrak m$. The Faddeev--Popov
ghosts that accompany the local $\mathrm{Spin}(10)\times U(1)$ gauge fixing
are Lorentz scalars in the adjoint of $H$, and their fluctuation operator is
the \emph{adjoint covariant Laplacian},
\begin{equation}\label{eq:gravity:ghost}
  \Delta_{\mathrm{ghost}} \;=\; -\bigl(D_{\mathrm{adj}}\bigr)^{2},
  \qquad
  D_{\mathrm{adj}} = d + [A_{\mathrm{adj}},\,\cdot\,] ,
\end{equation}
acting on $\mathrm{adj}(H)$. The ghosts are massless and contribute to the
subtracted quadratic divergence but, being massless, do \emph{not}
contribute to the curvature ($a_{2}$) running --- see
\autoref{thm:gravity:planck} and \autoref{app:gravity:main}; no quoted
heat-kernel number changes.

%----------------------------------------------------------------------%
\subsection{Induced Planck mass from the $a_{2}$ coefficient}
\label{sec:gravity:planck}
%----------------------------------------------------------------------%

We evaluate $\tfrac12\Tr\log\Delta_{m}$ in the proper-time (Schwinger)
representation on the slowly varying curved background $M^{4}$. The
heat-kernel expansion of a real scalar of mass $m$ is
\begin{equation}\label{eq:gravity:hk}
  \Tr\, e^{-s\Delta_{m}}
  \;=\;
  \frac{1}{(4\pi s)^{2}}\, e^{-s m^{2}}
  \int_{M^{4}}\!\sqrt{-g}\,
  \Bigl[\,a_{0} + a_{1}\,s + a_{2}\,s^{2} + \cdots\,\Bigr],
\end{equation}
with the Seeley--DeWitt coefficients (\autoref{app:gravity:main})
\begin{equation}\label{eq:gravity:a2}
  a_{0} = 1, \qquad
  a_{1} = \tfrac{1}{6}\,R, \qquad
  a_{2} = \tfrac{1}{180}\bigl(R_{\mu\nu\rho\sigma}^{2} - R_{\mu\nu}^{2}\bigr)
        + \tfrac{1}{72}R^{2} + \tfrac{1}{6}\,\Box\!\left(\tfrac{R}{5}\right)+\cdots
\end{equation}
The Einstein--Hilbert term comes from the curvature-linear coefficient
$a_{1}=\tfrac16 R$; its coefficient in the effective action is the
$n=1$ proper-time moment. Because the $a_{n}$ are
\emph{mass-independent} geometric invariants, the entire mass dependence
sits in the $e^{-sm^{2}}$ factor, and the curvature-linear term integrates
to a logarithm cut off at the scale $\mu$. Crucially this is finite and
non-zero \emph{only} for $m\neq 0$: a massless field gives $e^{-s m^{2}}=1$
and the $R$-coefficient reduces to a pure (subtracted) quadratic divergence
$\sim\Lambda^{2}$ with no $m^{2}\log$ running. Massless fields (Goldstones,
gauge bosons, ghosts) therefore do not contribute to $M_{\mathrm{Pl}}^{2}$.

\begin{theorem}[Induced Planck mass; Sakharov]
\label{thm:gravity:planck}
\label{thm:gravity:MPl}
\label{thm:GravitySummary}
Given the soldered background of Postulate~P3 (\autoref{sec:spacetime:main})
and the $28$ massive clock fields of \autoref{thm:gravity:spectrum}, the
one-loop effective action contains an Einstein--Hilbert term
\begin{equation}\label{eq:gravity:EH}
  \Gamma^{(1)} \;\supset\;
  \frac{M_{\mathrm{Pl}}^{2}}{2}\int_{M^{4}}\!\sqrt{-g}\;R,
\end{equation}
with induced Planck mass
\begin{equation}\label{eq:gravity:MPl}
  M_{\mathrm{Pl}}^{2}
  \;=\;
  \frac{1}{16\pi^{2}}\cdot\frac{1}{6}
  \sum_{\text{massive } i} s_{i}\, m_{i}^{2}\,\log\frac{\Lambda^{2}}{m_{i}^{2}}
  \;=\;
  \frac{28}{6}\,
  \frac{f^{2}\mu^{2}\,\log(\Lambda^{2}/\mu^{2})}{16\pi^{2}}
  \;>\; 0 ,
\end{equation}
where $s_{i}=+1$ for scalars, $\Lambda$ is the ultraviolet cutoff, and the
$28$ scalars all carry $m^{2}=f^{2}\mu^{2}$.
\end{theorem}

\begin{proof}
The curvature-linear term in $\tfrac12\Tr\log\Delta_{m}$ is, by
\eqref{eq:gravity:hk}--\eqref{eq:gravity:a2}, the $n=1$ proper-time moment
of $a_{1}=\tfrac16 R$. Writing the proper-time integral with the cutoff
$\Lambda$ in the lower limit and the mass scale in the renormalization
condition, the (a priori $0/0$-singular) $n=1$ term evaluates explicitly to
\begin{equation}\label{eq:gravity:n1}
  \bigl[\,n=1\ \text{term}\,\bigr]
  \;=\;
  a_{1}\cdot\frac12\log\frac{\Lambda^{2}}{\mu^{2}}
  \;=\;
  \frac{1}{6}\,R\cdot\frac12\log\frac{\Lambda^{2}}{\mu^{2}},
\end{equation}
per real scalar of mass $m^{2}=\mu^{2}$ in canonical normalization (the
$\tfrac16$ is the conformal Seeley--DeWitt coefficient). Multiplying by the
loop measure $1/(16\pi^{2})$ and summing $s_{i}=+1$ over the $28$ massive
clock fields, each with $m_{i}^{2}=f^{2}\mu^{2}$, gives the prefactor
$28/6$ and the common log $\log(\Lambda^{2}/\mu^{2})$, hence
\eqref{eq:gravity:MPl}. The detailed proper-time bookkeeping is in
\autoref{app:gravity:main}.
\end{proof}

\begin{proposition}[Positivity of the induced Newton constant]
\label{prop:gravity:positivity}
$M_{\mathrm{Pl}}^{2}>0$ for $\Lambda>\mu$.
\end{proposition}

\begin{proof}
Every factor in \eqref{eq:gravity:MPl} is strictly positive:
$\tfrac{28}{6}=4.6667>0$, $1/(16\pi^{2})=0.0063326>0$, $f^{2}\mu^{2}>0$
(since $[f]=\text{mass}$ and $\mu$ real), and
$\log(\Lambda^{2}/\mu^{2})>0$ for $\Lambda>\mu$. The $28$ scalars
($s=+1$) dominate any massive-fermion negative ($s=-1$) contribution, so the
induced Newton constant $G_{N}=1/(8\pi M_{\mathrm{Pl}}^{2})$ has the correct
attractive sign. A worked representative ($f=\mu=1$ in $\mu$-units,
$\Lambda/\mu=100$, so $\log(\Lambda^{2}/\mu^{2})=\log 10^{4}=9.2103$) gives
$M_{\mathrm{Pl}}^{2}=4.6667\times 9.2103\times 0.0063326=0.27218\,\mu^{2}>0$.
\end{proof}

\begin{remark}[Running of $M_{\mathrm{Pl}}^{2}$ from the same expression]
\label{rmk:gravity:running}
The static value and the running are derived from the \emph{one}
expression \eqref{eq:gravity:MPl}, so there is no $32$-versus-$28$ or
normalization mismatch. Differentiating with respect to the renormalization
scale,
\begin{equation}\label{eq:gravity:betaMPl}
  \beta_{M_{\mathrm{Pl}}^{2}}
  \;=\;
  \mu\,\frac{d M_{\mathrm{Pl}}^{2}}{d\mu}
  \;=\;
  -\,\frac{28}{6}\,\frac{2 f^{2}\mu^{2}}{16\pi^{2}}
  \;+\;\frac{28}{6}\,\frac{2 f^{2}\mu^{2}\log(\Lambda^{2}/\mu^{2})}{16\pi^{2}},
\end{equation}
so $M_{\mathrm{Pl}}^{2}(\mu)$ grows (logarithmically dominated) and remains
positive throughout the regime $\mu\ll\Lambda$. The prefactor $28/6$ is the
same number that appears statically.
\end{remark}

\begin{remark}[Why $28/6$ and not $32$]
\label{rmk:gravity:2832}
The legacy value used $N_{s}=32$ Goldstones and dropped the conformal
$\tfrac16$. That is wrong on two counts: only the $28$ \emph{physical
massive} clock fields run the $R$-coefficient (the $4$ frame modes are eaten
and massless; massless fields give no $m^{2}\log$), and the per-scalar
heat-kernel coefficient is $a_{1}=\tfrac16 R$, so the $\tfrac16$ and the
$\mu^{2}$ must be kept. The correct, manuscript-wide prefactor is
$28/6$.
\end{remark}

%----------------------------------------------------------------------%
\subsection{Cosmological constant: inherited, not solved}
\label{sec:gravity:cc}
%----------------------------------------------------------------------%

The curvature-independent coefficient $a_{0}=1$ in \eqref{eq:gravity:hk}
produces a vacuum-energy (cosmological-constant) term whose proper-time
$n=-1$ moment is quartically divergent,
\begin{equation}\label{eq:gravity:Lambda4}
  \Gamma^{(1)} \;\supset\;
  \int_{M^{4}}\!\sqrt{-g}\;\rho_{\mathrm{vac}},
  \qquad
  \rho_{\mathrm{vac}}
  \;\sim\;
  \frac{1}{16\pi^{2}}\sum_{i} s_{i}\,\Lambda^{4}
  \;\sim\;
  \frac{28}{32\pi^{2}}\,\Lambda^{4},
\end{equation}
with \emph{no} supersymmetric cancellation in UCFT. This $\sim\Lambda^{4}$
vacuum energy is not removed by any symmetry of the construction.

\begin{postulate}[Cosmological-constant counterterm]
\label{post:gravity:cc}
A counterterm $\Lambda_{\mathrm{cc}}$ is added to cancel
\eqref{eq:gravity:Lambda4} down to the observed value. In the
renormalization-group language $\Lambda_{\mathrm{cc}}$ is a tuned
\emph{relevant} coupling (canonical eigenvalue $-4$ of the stability matrix,
\autoref{sec:gravity:rg}).
\end{postulate}

\begin{remark}[Status]
\label{rmk:gravity:ccstatus}
UCFT \emph{inherits}, and does \emph{not} solve, the cosmological-constant
problem: the magnitude of $\Lambda_{\mathrm{cc}}$ is fixed by hand, exactly
as in the Standard Model coupled to gravity. No claim of naturalness is
made.
\end{remark}

%----------------------------------------------------------------------%
\subsection{Place in the renormalization-group flow}
\label{sec:gravity:rg}
\label{sec:gravity:stability}
%----------------------------------------------------------------------%

The gravitational coupling $\kappa$ (dimensionless Newton coupling) joins
the gauge--Yukawa flow. The gauge sector is asymptotically \emph{free}
(Gaussian fixed point $\hat g^{2}_{\star}=0$; full one- and two-loop
analysis in \autoref{sec:frg:main}); here we record only the gravitational
beta function and the combined fixed point, with the truncation caveat below.

\begin{conjecture}[Gravitational fixed point]
\label{conj:gravity:as}
\label{sec:gravity:as}
\label{sec:gravity:AS}
The single-scalar-loop gravitational beta function is
\begin{equation}\label{eq:gravity:betakappa}
  \beta_{\kappa} \;=\; 2\kappa - \frac{5}{48\pi^{2}}\,\kappa^{2},
\end{equation}
with a nontrivial ultraviolet-attractive fixed point
\begin{equation}\label{eq:gravity:kappastar}
  \kappa_{\star} \;=\; \frac{96\pi^{2}}{5} \;=\; 189.4964,
  \qquad
  \frac{d\beta_{\kappa}}{d\kappa}\Big|_{\kappa_{\star}}
  \;=\; 2 - \frac{2\cdot 5}{48\pi^{2}}\,\kappa_{\star}
  \;=\; -2 .
\end{equation}
The negative sign in \eqref{eq:gravity:betakappa} is essential: $\kappa$
flows \emph{away} from $0$ toward $\kappa_{\star}$ (gravity does \emph{not}
decouple), consistent with the anomalous dimension
$\eta_{M}=-(5/48\pi^{2})\kappa$.
\end{conjecture}

\begin{remark}[Caveat]
\label{rmk:gravity:caveat}
$\kappa_{\star}=96\pi^{2}/5\approx 189\gg 1$ lies outside the controlled
weak-coupling regime, so this fixed point is \emph{indicative}, not
rigorous. We label the gravitational asymptotic-safety statement a
conjecture for this reason.
\end{remark}

Collecting the couplings $(\xi,\hat g^{2},\hat y^{2},\kappa,\Lambda_{\mathrm{cc}})$
and using the asymptotic-safety convention $\theta_{i}=-\mathrm{eig}(M)$
(UV-attractive $\iff \mathrm{eig}(M)>0$; $\#\text{relevant}=\#\{\mathrm{eig}(M)<0\}$),
the representative fixed point self-consistent with the gauge asymptotic
freedom of \autoref{sec:frg:main} and with \eqref{eq:gravity:kappastar} is
\begin{equation}\label{eq:gravity:FP}
  (\xi_{\star},\,\hat g^{2}_{\star},\,\hat y^{2}_{\star},\,\kappa_{\star},\,\Lambda_{\mathrm{cc},\star})
  \;=\;
  \bigl(0.05,\ 0,\ 0,\ 96\pi^{2}/5,\ 0.01\bigr).
\end{equation}
At this point $\hat g^{2}_{\star}=\hat y^{2}_{\star}=0$ are Gaussian
(asymptotic freedom), so the stability matrix $M=\partial\beta_{i}/\partial g_{j}$
is effectively triangular with diagonal entries
\begin{equation}\label{eq:gravity:stab}
  \frac{\partial\beta_{\xi}}{\partial\xi}=+2,\quad
  \frac{\partial\beta_{\hat g^{2}}}{\partial\hat g^{2}}=+2,\quad
  \frac{\partial\beta_{\hat y^{2}}}{\partial\hat y^{2}}=+2,\quad
  \frac{\partial\beta_{\kappa}}{\partial\kappa}=-2,\quad
  \frac{\partial\beta_{\Lambda_{\mathrm{cc}}}}{\partial\Lambda_{\mathrm{cc}}}=-4,
\end{equation}
so that
\begin{equation}\label{eq:gravity:eig}
  \mathrm{eig}(M)=\{+2,+2,+2,-2,-4\},
  \qquad
  \theta=-\mathrm{eig}(M)=\{-2,-2,-2,+2,+4\}.
\end{equation}
The two relevant directions ($\#\{\mathrm{eig}(M)<0\}=2$) are $\kappa$
(eigenvalue $-2$, $\theta=+2$, the gravitational fixed point) and
$\Lambda_{\mathrm{cc}}$ (eigenvalue $-4$, $\theta=+4$, the tuned
cosmological-constant counterterm of \autoref{post:gravity:cc}). The
positive central $(\hat g^{2},\hat y^{2})$ eigenvalues are the correct
signature of asymptotic freedom under this convention. Ultraviolet
completeness of UCFT is therefore asymptotic freedom of the gauge--Yukawa
sector \emph{plus} the (indicative) gravitational fixed point \emph{plus}
irrelevant higher-curvature operators --- \emph{not} a non-Gaussian gauge
fixed point.

\begin{remark}[Higher-curvature operators]
\label{rmk:gravity:higher}
The $a_{2}$ coefficient \eqref{eq:gravity:a2} supplies curvature-squared
terms $c_{1}R^{2}+c_{2}R_{\mu\nu}R^{\mu\nu}$ with
$|c_{1,2}|\sim 1/(16\pi^{2})$, suppressed by $\mu^{2}/\Lambda^{2}$ relative
to $M_{\mathrm{Pl}}^{2}R$ and hence irrelevant in the regime
$R\ll M_{\mathrm{Pl}}^{2}$. These are the irrelevant operators of the
ultraviolet-completeness statement above.
\end{remark}
